% Independent ordinary-mathematical audit; meant for inclusion in report.tex.
% No new proof-assistant build is asserted.
\section{Logical endpoint and arbitrary-root phase erasure}
\label{sec:logical-audit}

\subsection{Scope of the inherited consistency theorem}

Let $C_{254}$ be the original numerical calculus and let $F=1$ be its
designated refutation code.  The source program retains truth, modus ponens,
generalization, all eight logical rows, the native equality mechanism,
both freshness tests, and the complete CP/GIC bundle.  The native theory
has the CP container axiom, every formula instance of global-image
collection, and raw irreflexivity.  In particular, GIC ranges over finite
syntactic formulas; it is not collection for every external binary relation.

The source audit verifies that the inherited composition has the form
\[
 \operatorname{Halt}_{\mathrm{source}}(R)
 \ \longleftrightarrow\ C_{254}(1)
 \ \longleftrightarrow\ \operatorname{Refute}(\mathrm{NativeAx})
 \ \longleftrightarrow\ \operatorname{Refute}(\mathrm{ZF})
 \ \longleftrightarrow\ \operatorname{Refute}(\mathrm{ZFC}),
\]
provided the initial source register data are safe.  The first equivalence
uses the actual path program and the literal-to-structured program
comparison.  The next two use the native-calculus comparison and the
internally well-founded part of the bisimulation quotient.  The last uses
the constructible subdomain of an arbitrary set-coded ZF model.

These model-theoretic arguments are available inside ZFC.  Satisfaction for
a set structure is defined by recursion on finite syntax; the quotient and
the definable subdomains are sets of the ambient metatheory.  Ordinary
first-order completeness supplies a nonempty model when a recursively
axiomatized countable theory is consistent.  Neither external
well-foundedness nor transitivity of that model is required.  The
well-foundedness predicate used inside the quotient quantifies over the
subsets represented by that model, and its closure lemmas use the actual
GIC instances.  The constructible subdomain is its internal $L$, with its
own set-theoretic operations and Choice theorem.

This is an ordinary ZFC proof route for the displayed equivalence.  The
external Lean declarations identify the mathematics being composed; their
standard axiom reports alone would not constitute a general translation
from Lean into ZFC.  A PA proof additionally needs represented proof-code
and execution translations.  The source explicitly leaves that stronger
arithmetization task open.

\subsection{The exact safe-root specialization}

Use logical register order
\[
 (N,T,W,A,D,H,X,Y,L).
\]
The predicate used by \texttt{DirtyProgram254.safe\_iff} is exactly
\[
 C_{254}(T)\ \wedge\ C_{254}(H)\ \wedge\operatorname{Good}(W).
\]
All six remaining fields may be arbitrary natural numbers.  The current
246-state bootstrap has logical root
\[
 (3,0,0,0,3,0,0,0,0).
\]
It therefore instantiates
\texttt{six\_nonformula\_path\_iff} at
$(N,A,D,X,Y,L)=(3,0,3,0,0,0)$, with two truth constructors and
\texttt{Enum254.good\_zero}.  Its physical register permutation is
$(6,2,0,4,3,7,8,5,1)$; the scalar nonzero value in physical block three
is the numerical parameter $D$, not a formula or stack field.

The proof permits arbitrary scratch values by replacing a dirty root with
an explicitly computed virtual clean root that reaches the same state
after the fetch and first read.  This map preserves $T,H,W$.  Completeness
from an arbitrary candidate counter follows by padding an initial reset
command: a refutation has arbitrarily large valid command encodings, and
the scheduler reaches any prescribed later candidate unless it has
already halted.  Thus starting at $N=3$ loses no refutation.

The logical register file above index eight is zero in \texttt{L2.enc}.
Physical spare blocks are a different representation resource and may
have the values allowed by the binary frame contracts.  These two tails
must not be identified when composing the source and tape semantics.

\subsection{A strengthened phase invariant}

The earlier phase analysis identifies two transfer loops whose positive
guards never occur from the zero-register root.  The present calculation
strengthens this to \emph{every} initial nine-register vector at the root.
This removes the zero-initialization qualification from the source-level
optimization, while retaining the distinction between root starts and
arbitrary interior starts.

For each register, use the finite abstract domain consisting of unions of
the three classes $\{0\}$, $\{1\}$, and $\{2,3,\ldots\}$.  A program
point holds a Cartesian product of nine such unions, or the empty set.
Increment sends $0$ to $1$ and positive values to the third class.
Positive decrement sends $1$ to $0$ and the third class to the union of
the second and third classes.  Zero decrement filters to $0$ and changes
no register.  Break and continue change only the finite control.

Each abstract transfer contains every concrete successor of every value
in its input set.  Componentwise union is a sound join.  Start the worklist
with the full nine-fold product $\mathbb N^9$ at the leaf decoded from the
empty PC.  The worklist can add only finitely many class bits and therefore
terminates.  More significantly, the final certificate is checked directly:
it contains the initial product, and the image of every feasible guarded
edge lies inside the mask at its actual successor.  Induction on source
steps proves that every concrete root execution lies in these masks.

The all-registers-arbitrary computation terminates after 35,768 worklist
pops.  Its postfix certificate admits 6,456 of the 6,460 leaves and 8,385
of the 8,391 guards.  Its only exclusions remain
\[
 8202,\ 8203,\ 8204,\ 8358,\ 8359,\ 8360.
\]
The two tested registers are exactly zero, independently of every other
register value.  This is an unbounded invariant: there is no bound on
register magnitudes or on the number of computation steps.

\begin{center}
\begin{tabular}{rlll}
Old node & Loop & Impossible positive case & Zero continuation\\
\hline
6324 & $X\mathrel{+\!\to}H$ & 8202 & increment $H$\\
6436 & $T\mathrel{+\!\to}X$ & 8358 & begin transfer $W\to T$
\end{tabular}
\end{center}

The notation $U\mathrel{+\!\to}V$ means consume $U$ while adding its
old value to $V$.  Both loops are transfers, rather than destructive
assignments that could discard an arbitrary nonzero source.  Because their
source is zero at every root-reachable call, their actual effect is the
identity on the complete register vector.

\subsection{The concrete occurrence-level contraction}

In the literal \texttt{Prog254.v} tree, the relevant fragments are
\[
\begin{array}{rcl}
6323&=&\operatorname{seq}(6324,\operatorname{inc}5),\\
6324&=&\operatorname{seq}(\operatorname{decnz}6,6325),\\
6325&=&\operatorname{seq}(\operatorname{inc}5,\operatorname{continue}2),\\[2pt]
6435&=&\operatorname{seq}(6436,6438),\\
6436&=&\operatorname{seq}(\operatorname{decnz}1,6437),\\
6437&=&\operatorname{seq}(\operatorname{inc}6,\operatorname{continue}2).
\end{array}
\]
Replace occurrence 6323 by $\operatorname{inc}5$, and occurrence 6435 by
subtree 6438.  The qualifier ``occurrence'' matters: other uses of the
shared transfer macros are unchanged and may have nonzero inputs.

The generated full tree has 6,453 internal nodes, 6,454 instruction leaves,
and 8,383 guarded cases.  Exactly six old leaves disappear: two always-zero
decrement guards and the four unreachable body leaves.  Existing break
and continue counts are retained.  The generator checks all 8,383 remaining
guarded control successors after rebasing every retained leaf path.  They
agree exactly with the old successors after bypassing the two zero guards.
The check includes the final carry zeros in the PC representation and all
break and continue operations; no trailing-zero truncation is substituted.

\paragraph{Erasure theorem.}
Let $P$ be the original literal path program and let $P^{\dagger}$ be this
generated tree.  For every nine-register vector $R$, both programs started
at the empty instruction pointer halt alike:
\[
 \operatorname{Halt}(P,[],R)
 \quad\longleftrightarrow\quad
 \operatorname{Halt}(P^{\dagger},[],R).
\]

\paragraph{Proof.}
The postfix certificate gives the preceding phase invariant on every
original root run.  At either deleted decrement, its source register is
zero, so the next state has identical registers and follows the certified
zero successor.  The two deleted positive branches and their bodies are
unreachable.  Each retained instruction has the same operation and register
argument in the new tree.  The verified path correspondence identifies its
new successor with the original successor after these identity steps.
This proves forward simulation by induction along a finite run.

For the reverse direction, start with the same initial register vector.
The successor certificate and the concrete transfer laws expand every new
step into the corresponding old step, followed by zero or one deleted
identity step.  The transported invariant supplies zero in the tested
register whenever such a step is inserted.  Induction constructs an old
run with the same registers at every retained point.  There is at most one
erased step between consecutive retained points, and both programs retain
the same unique terminal leaf.  Consequently the expansion or contraction
of a finite halting run is finite and halting.  This proves the equivalence.
\hfill$\square$

The theorem applies without formula-safety assumptions; those assumptions
are needed only when its common halting behavior is calibrated to
$C_{254}(1)$ and hence to ZFC inconsistency.  The proof consists of finite
control equalities and elementary induction, so it supplies an ordinary
ZFC argument and a concrete proof-assistant port target.

\subsection{Compiler measurement and the remaining question}

Shorter source trees need not give smaller compressed programs.  Exact
subtree sharing of the original literal tree yields 203 distinct pair
subtrees.  Contracting only the $X\to H$ occurrence yields 206; contracting
only the $T\to X$ occurrence yields 204; contracting both yields 207.
In the last comparison, fifteen old unique pairs disappear but nineteen
context-specific pairs are created.  The transfer macros still occur
elsewhere.  Thus plain subtree hashing increases the decision graph by
four nodes in this particular transformation.

This negative measurement is useful: the sound source reduction is
available, but obtaining a binary-state saving requires further grammar or
layout optimization.  The delivered tree and certificates do not by
themselves claim a new binary Turing-machine state count, a new Lean/Rocq
endpoint, or a proof object internal to PA.  The source-level erasure is
kept separate from any successful binary optimization found elsewhere in
this report.

\subsection{Independence hypotheses}

Suppose ZFC proves
$\neg\operatorname{Halt}(M)\leftrightarrow\operatorname{Con}(\mathrm{ZFC})$,
where consistency is expressed using the standard proof predicate.
Under the usual hypotheses of the second incompleteness theorem, external
consistency of ZFC then prevents ZFC from proving
$\neg\operatorname{Halt}(M)$: such a proof, together with the internal
equivalence, would give a ZFC proof of its own consistency.  This is a
syntactic unprovability conclusion.

Actual nonhalting uses a separate external correctness argument.  For the
proof-search construction discussed here, the external calibration shows
that a finite halting run yields a finite ZFC contradiction proof.
External consistency therefore rules out such a run.  The bare internal
equivalence and external consistency alone do not establish this semantic
conclusion.  The calibration invoked here is the ordinary simulation
argument already described, not a new proof-code arithmetization.

Once actual nonhalting is established, $\Sigma_1$-soundness of ZFC excludes
a ZFC proof of the false $\Sigma_1$ statement that $M$ halts.  Consistency
alone does not exclude such a false existential theorem.  Thus a full
two-sided independence claim must state its stronger hypothesis
explicitly; the internal equivalence theorem itself requires no
soundness assumption.
